\section{Reciprocal remainders and Gaussian odd parts}
\label{rlc:sec:harmonic}

\subsection{An all-order identity for polylogarithmic remainders}
For $N\ge0$ define the finite Taylor remainder
\begin{equation}\label{rlc:eq:remainderdef}
 \mathcal R_{N,s}(y)=\Li_s(-y)-\sum_{k=1}^{N}\frac{(-y)^k}{k^s}.
\end{equation}
The empty sum at $N=0$ is zero. The elementary Lerch shift formula
implies
\begin{equation}\label{rlc:eq:remaindershift}
 F_{N+1,s}(y)=-(-y)^{-N}\mathcal R_{N,s}(y).
\end{equation}
Although each remainder can grow polynomially on the positive real
axis, the reciprocal product is integrable.

\begin{corollary}[Harmonic-tail product identity]
\label{rlc:cor:remainders}
For every $N\ge0$ and $s,t\in\C$,
\begin{equation}\label{rlc:eq:remainderidentity}
 \boxed{\displaystyle
 \int_0^\infty\mathcal R_{N,s}(y)\mathcal R_{N,t}(1/y)
                  \frac{\dd y}{y}
 =\E_{N+1}(s+t)
 =(s+t)\left[\zeta(s+t+1)-H_N^{(s+t+1)}\right].}
\end{equation}
The last expression is filled in by its removable value $1$ when
$s+t=0$. Every order derivative and every logarithmic moment is obtained
by substituting $a=N+1$ into the preceding theorems.
\end{corollary}
\begin{proof}
The two powers $(-y)^{-N}$ and $(-1/y)^{-N}$ in
\eqref{rlc:eq:remaindershift} multiply to one. Apply
\eqref{rlc:eq:opening}, followed by the Hurwitz shift identity
$\zeta(v,N+1)=\zeta(v)-H_N^{(v)}$. Normal convergence already proved
for the Lerch profiles gives all differentiated versions.
\end{proof}

At $N=1$ and $s=t=1$, this gives the elementary-looking but nontrivial
identity
\begin{equation}\label{rlc:eq:remainderexample}
 \int_0^\infty
 [\log(1+y)-y][\log(1+1/y)-1/y]\frac{\dd y}{y}
       =2\zeta(3)-2.
\end{equation}
The factors must remain combined. Expanding their product and
integrating the resulting pieces independently introduces divergent
integrals and is not a valid proof of the formula.

The Stieltjes constants at the shifted integer arguments also reduce
by a finite harmonic-logarithmic correction:
\begin{equation}\label{rlc:eq:stieltjesshift}
 \gamma_n(N+1)=\gamma_n(1)
           -\sum_{k=1}^N\frac{\log^n k}{k}.
\end{equation}
It follows by comparing Laurent coefficients of
$\zeta(1+v,N+1)=\zeta(1+v)-\sum_{k=1}^N k^{-1-v}$.
For $n=0$ it reads $\gamma_0(N+1)=\gamma-H_N$.
Thus the entire remainder family contains finite generalized harmonic
numbers and their order derivatives without changing the convergence
normalization.

\subsection{A fourth-root polylogarithm family}
For complex $s$ and positive $y$, define the analytic odd part
\begin{equation}\label{rlc:eq:Adef}
 A_s(y)=\frac{\Li_s(iy)-\Li_s(-iy)}{2i}.
\end{equation}
For real $s$ this agrees with an imaginary part. For complex $s$,
\eqref{rlc:eq:Adef}, rather than $\Ima\Li_s(iy)$, is required for
holomorphic order differentiation.
The initial series gives
\[
 A_s(y)=\sum_{n=0}^\infty\frac{(-1)^ny^{2n+1}}{(2n+1)^s},\qquad
 F_{1/2,s}(y^2)=2^s y A_s(y),
\]
with continuation to the positive real axis.

\begin{corollary}[Gaussian reciprocal order addition]
\label{rlc:cor:gaussian}
For all $s,t\in\C$, with $W=s+t$,
\begin{equation}\label{rlc:eq:gaussian}
 \boxed{\displaystyle
 \int_0^\infty A_s(y)A_t(1/y)\frac{\dd y}{y}
 =2^{-W-1}\E_{1/2}(W)
       =(1-2^{-W-1})\E_1(W).}
\end{equation}
Every mixed order derivative is an ordinary convergent integral.
\end{corollary}
\begin{proof}
Use the displayed relation to $F_{1/2,s}$ and set $z=y^2$.
The measure contributes $1/2$ and the order prefactors contribute
$2^{-W}$. Theorem~\ref{rlc:thm:addition} at depth two gives the first
equality. The second follows from
$\zeta(v,1/2)=(2^v-1)\zeta(v)$, with the pole cancelled before evaluation.
\end{proof}

Since $A_1(y)=\arctan y$, a basic specialization is
\begin{equation}\label{rlc:eq:arctan}
 \int_0^\infty\arctan y\,\arctan(1/y)\frac{\dd y}{y}
     =\frac74\zeta(3).
\end{equation}
More characteristic of the order-derivative continuation is the identity
\begin{equation}\label{rlc:eq:gaussianjets}
 \int_0^\infty A_0^{[1]}(y)A_0^{[1]}(1/y)\frac{\dd y}{y}
 =-\gamma_1+\gamma\log2-\frac12\log^22.
\end{equation}
Here $A_0^{[1]}=\left.\partial_sA_s\right|_{s=0}$.
It follows by taking the second derivative at zero of
$(1-2^{-W-1})\E_1(W)$.

These identities use fourth-root arguments, but their relation to the
repository's Gaussian depth-two periods is not automatic. In particular,
\eqref{rlc:eq:gaussian} is not a reduction of
$\sum_{n\ge0}(-1)^nH_n/(2n+1)^6$ in the prescribed $g_{a,b}$ basket.
That distinction is retained in the integration status.
